% module_id: transport.dd
% module_version: 0.1.0
% plugin_dir: plugins/transport_dd
% source_files: plugins/transport_dd/{__init__,plugin,physics,assemble,newton}.py
% tests: tests/test_transport_dd.py

\chapter{稳态漂移扩散输运}
\label{ch:transport_dd}

\begin{quote}
\small
\noindent\textbf{模块标识}\quad\pid{transport.dd}\qquad
\textbf{版本}\quad 0.1.0\qquad
\textbf{类别}\quad transport
\end{quote}

\noindent
本章记述电荷\emph{如何流动}。前面的极化与量子模块回答了"电荷从哪来"，
给出一个平衡态的电荷分布；本模块把它扩展成\emph{带偏压}的问题：
在结构化二维网格上联立求解静电 Poisson 方程与电子连续性方程，
扫过栅、漏偏压，产出电流--电压曲线。

这是全流水线中数值上最重的一章。它的未知量数是前几章的
$O(10^{4})$ 倍，Jacobian 是块三对角的，求解过程需要三层递进的
保险机制才不至于发散。相应的，它的验证也最系统：六个级别的基准
（T0--T5）全部先于任何"像不像器件"的检查。

\section{模块契约}
\label{sec:dd-contract}
\index{4module@模块!transport.dd}

\begin{modulecard}{transport.dd}{0.1.0}
  \cardrow{类别}{transport（输运求解，流水线阶段）}
  \cardrow{能力}{transport, transport.dd, current, iv}
  \cardrow{依赖 artifact}{mesh, structure, materials, polarization}
  \cardrow{产出 artifact}{electrostatics, carriers, current, iv}
  \cardrow{验证状态}{partially-verified}
  \cardrow{源码}{\fpath{plugins/transport_dd/}（多文件插件包）}
  \cardrow{测试}{\fpath{tests/test_transport_dd.py}}
  \cardrow{可标定参数}{\akey{temperature_K}\calib、\akey{mu0_cm2_Vs}\calib、
    \akey{vsat_cm_s}\calib、\akey{delta_ec_ratio}\calib}
  \cardrow{协作者}{经 Deck 的 \akey{linear} 槽解析 \pid{linear.blocktridiag}
    （\crefchap{ch:blocktridiag}）}
  \cardrow{模型范围}{稳态、等温、单极（仅电子）}
\end{modulecard}

本模块为何是 \texttt{partially-verified} 而不是 \texttt{unverified}，
以及为何还不是 \texttt{verified}，模块说明书自己给出了判据：
设计稿原本的门槛是"T0--T4 全部通过之前保持 unverified"，它们确实通过了
（均匀电阻闭式解、制造解二阶收敛、端点电流连续性、延拓无关性）；
而阻止它升级为 \texttt{verified} 的，是\emph{数值}所依赖的迁移率未经标定、
\texttt{statistics} 只有 Boltzmann，以及两条设计判据在本离散下
\emph{原理上不可达}（见 \S\ref{sec:dd-limits}）。

\section{控制方程与物理假设}
\label{sec:dd-model}

\subsection{方程组}
\label{sec:dd-equations}
\index{1concept@概念!drift-diffusion@漂移扩散}

在稳态、等温、单极（只考虑电子）的假定下，器件由两个方程描述：

\begin{equation}
\nabla\cdot\bigl(\epsr\epsO\nabla\psi\bigr) \;=\; -q\,\bigl(N_{\text{net}} - n\bigr),
\label{eq:dd-poisson}
\end{equation}

\begin{equation}
\nabla\cdot\Jnn \;=\; 0,
\qquad
\Jnn \;=\; q\,\mumu\,n\,\mathbf{E} \;+\; q\,D_n\nabla n ,
\label{eq:dd-continuity}
\end{equation}

$\psi$ 为静电势，$N_{\text{net}}=N_D^{+}-N_A^{-}$ 为净电离杂质浓度，
$\mathbf{E}=-\nabla\psi$。第二式（电流密度）在本模块中并不直接使用 ——
实际离散的是它的\emph{通量形式}，见 \S\ref{sec:dd-sg}。

"单极"这一限定是有代价的，且代价不小：电子单极模型里没有空穴去耗尽受主。
把缓冲层的 $1\times10^{17}$\U{cm^{-3}} 受主当作固定电离电荷，会在缓冲层内
造出非物理的巨大电场，而 Poisson 解隐含的电子密度又会让"半绝缘"缓冲层
\emph{变成导体} —— 实测表现为端点电流比沟道电流大一个数量级，
电流直接漏进衬底。因此本模块采用单极 HEMT 的标准抽象：
\textbf{缓冲层 = 中性介质 + 载流子钉零 + 对连续性方程的零通量壁}。

\subsection{为什么未知量取 $(\psi,n)$ 而不是 $(\psi,\varphi_n)$}
\label{sec:dd-unknowns}
\index{1concept@概念!unknown choice@未知量选择}

教科书通常以 $(\psi,\varphi_n)$（$\varphi_n$ 为准费米势）为未知量，
因为此时
$n=N_C\exp\!\bigl(q(\psi-\varphi_n)/kT\bigr)$ 使密度随 $\varphi_n$ 指数变化，
数值动态范围被压缩。本模块反其道而行，以\emph{密度本身}为未知量。理由在
\S\ref{sec:dd-sg} 的偏导结构里，此处先给出结论：

宽禁带势垒层与耗尽缓冲层中，密度会下溢到零。若以 $\varphi_n$ 为未知量，
连续性方程在这些区域会变成\emph{恒等于零}的空方程 —— 方程没坏，
但它不含任何信息，Newton 系统在这些行上是奇异的。而以 $n$ 为未知量时，
Scharfetter--Gummel 通量对密度的偏导\emph{与密度无关}
（\crefeq{eq:dd-flux-partials}），连续性方程在任何密度下都保持良态。

代价是未知量的动态范围极大，这由行缩放（\S\ref{sec:dd-scaling}）处理。

\subsection{统计}
\label{sec:dd-statistics}

载流子密度取 Boltzmann 近似：

\begin{equation}
n \;=\; \NC\,\exp\!\left(\frac{\psi-\varphi_n}{kT/q}\right),
\qquad
\NC \;=\; 2\left(\frac{2\pi \meff kT}{h^{2}}\right)^{3/2},
\label{eq:dd-boltzmann}
\end{equation}

$kT/q$ 以伏特计，使指数无量纲。指数在 $\pm100$ 处被截断；
这不是数值上的偷懒，而是一个\emph{可微}的钳制：截断之上实现残差
严格是"指数取截断值"的形式，因此 Jacobian 仍是它的精确导数，
Newton 不会因钳制而失去一致性。

参数 \akey{statistics} \textbf{只接受} \akey{boltzmann}，其余取值
显式报错。这是一个刻意留下的缺口：简并修正（Joyce--Dixon 或
Fermi--Dirac 积分）\emph{必须同时修改 Scharfetter--Gummel 通量}
（广义 Einstein 关系），只改密度公式会得到一个静默不一致的格式。
宁可拒绝，不要半搬。

\subsection{边界条件}
\label{sec:dd-boundary}

\begin{center}
\small
\begin{tabular}{@{}p{0.24\textwidth} p{0.66\textwidth}@{}}
\toprule
\textbf{边界} & \textbf{条件}\\
\midrule
源、漏接触 & 对 $\psi$ \emph{与} $n$ 同时施加 Dirichlet；两者随端电压同步平移，
  使接触面不携带伪电流\\
栅极 & 对 $\psi$ 的 Schottky Dirichlet（$v_{0.1.0}$ 无栅电流），
  对 $n$ 不加约束\\
背面 & 接地（$\psi=0$），衬底视为理想接地\\
无栅表面 & 钉在 Schottky 参考电位（工程近似，见 \S\ref{sec:dd-limits}）\\
缓冲层 & 载流子钉零，且对连续性方程是零通量壁\\
\bottomrule
\end{tabular}
\end{center}

"接触同时钉 $\psi$ 与 $n$"这一点值得强调。若只钉 $\psi$，接触面的
SG 通量会因两侧密度不匹配而产生一个虚假的电流，端点电流随即失去守恒性；
本模块的做法是把平衡态解在接触列上的取值\emph{整个}作为储层，
漏极在此基础上叠加 $V_{ds}$，因此接触面天然无电流。

\section{离散化}
\label{sec:dd-discretisation}

\subsection{有限体积五点格式}
\label{sec:dd-fv}
\index{3method@方法!finite volume@有限体积}

两个方程都在\emph{控制体}上积分。以 $V_i$ 记节点 $i$ 的控制体体积，
$A_{ij}$、$h_{ij}$ 记节点 $i$ 与其邻居 $j$ 之间面的面积与间距，则
Poisson 方程为

\begin{equation}
\sum_{j\in\mathcal{N}(i)}
\frac{\epsface_{ij}\epsO\,A_{ij}}{h_{ij}}
\bigl(\psi_i-\psi_j\bigr)
\;=\;
q\bigl(N_{\text{net},i}-n_i\bigr)V_i \;+\; \sigma_{\text{sheet}}\,\delta_{i,\text{int}},
\label{eq:dd-poisson-fv}
\end{equation}

面上的介电常数取两侧的\emph{调和}（串联）平均：

\begin{equation}
\epsface_{ij} \;=\; \frac{2\,\epsri{i}\,\epsri{j}}{\epsri{i}+\epsri{j}}.
\label{eq:dd-harmonic}
\end{equation}

调和平均是介电跃变与节点\emph{之间}时的正确处理；若界面恰在节点上，
它退化为算术平均。**本模块的界面恰在节点上**（\crefchap{ch:mesh}），
因此这一处在界面上被抹平一个格子宽度 —— 这是设计稿明示接受的代价，
其量化后果见 \S\ref{sec:dd-limits}。

连续性方程按通量守恒写：

\begin{equation}
\sum_{j\in\mathcal{N}(i)} J_{ij}\,A_{ij} \;=\; 0,
\label{eq:dd-continuity-fv}
\end{equation}

$J_{ij}$ 为由 Scharfetter--Gummel 公式给出的面电流密度。

\subsection{Scharfetter--Gummel 通量}
\label{sec:dd-sg}
\index{3method@方法!Scharfetter--Gummel 通量}

面电流不取 $\mathbf{E}$ 与 $\nabla n$ 的直接差分离散（那在大 Peclet 数下
会振荡），而取指数拟合形式 \cite{scharfetter1969}。以 $L$、$R$ 记面两侧：

\begin{equation}
\boxed{\;
J_{LR} \;=\; \frac{q\,\mumu\,kT}{h}
\Bigl[\,n_R\,B(\delta) \;-\; n_L\,B(-\delta)\,\Bigr],
\qquad
\delta \;=\; \frac{\Ec^{L}-\Ec^{R}}{kT},
\;}
\label{eq:sg-flux}
\end{equation}

\begin{equation}
B(x) \;=\; \frac{x}{\ee^{x}-1},
\qquad
B(-x) \;=\; B(x) + x .
\label{eq:bernoulli}
\end{equation}

\crefeq{eq:bernoulli} 中的\emph{反射恒等式}是 SG 权重的基本性质，也是本模块
踩过的一个坑：把反射关系写反（写成 $-B(x)+x$）之后，通量仍然
\emph{完全自洽} —— 系统级 Jacobian 对中心差分照样通过 $4\times10^{-10}$ ——
却在"带边沿电流方向上升"的每个面上系统性错误。修正后平衡态残差从 $35$
降到 $4.8\times10^{-6}$（七个数量级）。教训是：
\textbf{一致性不等于正确性}；通量类子程序必须对闭式解单独测试，
不能只靠"系统自洽"。

实现同时使用两个分支的解析导数 $B'(x)$，并且只对\emph{衰减}的指数求值：

\begin{equation}
B'(x) = \frac{\ee^{x}-1-x\ee^{x}}{(\ee^{x}-1)^{2}},
\qquad
B'(-x) = -B'(x)-1 \quad (x>0),
\label{eq:bernoulli-slope}
\end{equation}

$x\to0$ 处退回 Taylor 展开，避开 $0/0$。

\subsection{通量的解析偏导}
\label{sec:dd-flux-partials}

\crefeq{eq:dd-flux-partials} 列出通量对四个自变量的解析偏导，
它们是 Jacobian 的构件，也都被单独对中心差分验证：

\begin{equation}
\begin{aligned}
\pder{J_{LR}}{n_R} &= \frac{q\mumu kT}{h}\,B(\delta), &
\pder{J_{LR}}{n_L} &= -\frac{q\mumu kT}{h}\,B(-\delta),\\[2pt]
\pder{J_{LR}}{\psi_L} &= \frac{q\mumu kT}{h}\cdot\frac{1}{kT}
  \Bigl[n_R B'(\delta) + n_L B'(-\delta)\Bigr], &
\pder{J_{LR}}{\psi_R} &= -\,\pder{J_{LR}}{\psi_L}.
\end{aligned}
\label{eq:dd-flux-partials}
\end{equation}

三条性质值得逐条点明。

\textbf{（i）密度偏导与密度无关。} $\partial J/\partial n$ 只由几何、
迁移率与带边差决定，不含 $n$。这正是 \S\ref{sec:dd-unknowns} 选择 $n$
作未知量的理由，也解释了为什么宽禁带势垒或耗尽缓冲层里的连续性方程
仍然良态。

\textbf{（ii）势偏导反对称。} 把两个电位同时平移一个常数不改变 $\delta$，
也就不改变电流，因此 $\partial J/\partial\psi_L = -\partial J/\partial\psi_R$。
这个不变量是一个强判据：实现里只写了"带 $n_L$ 的那一项"的写法
是 $100\%$ 错误的，而它的密度偏导仍然精确 —— 正是那种"只在参考状态
做一次系统级 Jacobian 检查"会漏掉的错误。

\textbf{（iii）带边包含异质结台阶。} $\delta$ 中的 $\Ec$ 是
$\Ec^0-\psi$（数值上），因此 SG 通量\emph{同时}携带静电势差与
异质结带阶；平衡态下 $\delta$ 处处为零，通量恒等于零，
这是 T1 基准断言的不变量。

\subsection{场依赖迁移率}
\label{sec:dd-mobility}
\index{3method@方法!Caughey--Thomas 迁移率}

低场迁移率之外，本模块默认启用 Caughey--Thomas 速度饱和模型
\cite{caugheyThomas1967}：

\begin{equation}
\mumu(E) \;=\; \mumu_0\left[1+\left(\frac{\mumu_0 E}{v_{\text{sat}}}\right)^{\beta}\right]^{-1/\beta},
\label{eq:dd-mobility}
\end{equation}

低场下趋于 $\mumu_0$，高场下趋于 $v_{\text{sat}}/E$。
面电场取 $E=\abs{\Delta\Ec}/h$，即\emph{该面自己的}场。

关键在于它的导数必须进 Jacobian：SG 前因子带 $\mumu$，而面电场依赖带边，
所以 \crefeq{eq:dd-flux-partials} 的势偏导要多出一项

\begin{equation}
\pder{J}{\psi_L}\bigg|_{\mumu} \;=\;
\frac{q\,kT}{h}\Bigl[n_R B(\delta)-n_L B(-\delta)\Bigr]\pder{\mumu}{E}\pder{E}{\psi_L},
\qquad
\pder{E}{\psi_L} = -\frac{\operatorname{sgn}(\delta)}{h},
\label{eq:dd-mobility-chain}
\end{equation}

$\dd\mumu/\dd E$ 有闭式（在 $E=0$ 处保持有限，$\beta\ge1$）。
这一项是\emph{反对称}的，因此性质（ii）在启用场依赖后依然成立 ——
测试对此有单独断言，且**分别在启用与关闭场依赖两种情形下**做了
通量偏导的中心差分校验。

\subsection{块三对角结构与未知量排序}
\label{sec:dd-block}
\index{1concept@概念!block tridiagonal@块三对角}

每个节点携带两个未知量，按 $(\psi,\varphi_n)$ 交错存储；
竖直行 $j$ 展平后是 $2n_x$ 个未知量。张量积网格上，行 $j$ 只与
$j\pm1$ 耦合，因此 Jacobian 是\emph{块三对角}的，块大小 $m=2n_x$：

\begin{equation}
\mat{J} \;=\; \operatorname{tridiag}\bigl(\mat{L}_j,\ \mat{D}_j,\ \mat{U}_j\bigr),
\qquad \mat{L}_j,\mat{D}_j,\mat{U}_j \in \mathbb{R}^{m\times m}.
\label{eq:dd-block-structure}
\end{equation}

对演示网格（$n_x=41$，$n_z=1621$），未知量总数 $2n_xn_z = 132\,922$ ——
稠密分解不可行，而块三对角直接解的代价是 $O(n_zm^{3})=O(n_z(2n_x)^3)$，
这正是 \crefchap{ch:blocktridiag} 存在的理由。线性步经 Deck 的
\akey{linear} 槽解析，因此本模块并不知道、也不依赖具体用什么分解。

\subsection{行缩放：量纲化而非对角均衡}
\label{sec:dd-scaling}
\index{3method@方法!row scaling@行缩放}

两个方程块的单位不同（每单位宽度的电荷 vs.\ 每单位宽度的电流），
量级相差许多个数量级，直接取残差范数会失去意义。本模块对每个
非 Dirichlet 行除以\emph{固定的物理参考量}：

\begin{equation}
\text{Poisson 行} \div C_{\text{local}},
\qquad
\text{连续性行} \div I_{\text{ref}},
\label{eq:dd-row-scaling}
\end{equation}

前者把 Poisson 残差变成\emph{伏特}，后者把连续性残差变成\emph{相对电流}。
这是一次左乘正对角矩阵，不改变解，但使残差范数成为两个块共同的、
有物理意义的收敛度量。

有一条被\emph{否决}的做法值得记录：用 Jacobian 对角做行均衡
（即自适应缩放）。它会让宽禁带行与耗尽行获得巨大权重，
残差范数随即被相消噪声主导，线搜索于是接受"物理残差变差十二个数量级"的步。
实测把平衡态残差从 $4.2$ 改到 $1.3\times10^{-3}$ 才回到正轨。
未知量的\emph{列}缩放也被试过并否决，理由写在代码里：
在 $(\psi,n)$ 表述下矩阵元素与未知量本身无关
（$\dd J/\dd n$ 不含 $n$，Poisson 的 Laplacian 不含 $\psi$），
列缩放无法改善条件数，盲目施加反而会让块消元的枢轴溢出。

\section{求解算法}
\label{sec:dd-solver}

\subsection{三层保险}
\label{sec:dd-three-layers}
\index{3method@方法!pseudo-transient continuation@伪瞬态延拓}

方案把"Newton 收敛困难"列为最高风险，落地为三层递进的保险：

\begin{enumerate}[leftmargin=1.9em,itemsep=0.4em,topsep=0.4em]
  \item \textbf{偏置扫描延拓}：漏压（外层）与栅压（内层）逐点推进，
        前一点的收敛解作为下一点的初值；
  \item \textbf{阻尼线搜索}：残差不降则步长减半（Armijo 式接受），
        最多 \akey{damping_max_halvings} 次；
  \item \textbf{伪瞬态延拓（PTC）}：当线搜索连续失败时，在\emph{迭代矩阵}
        上注入物理时间项
        \begin{equation}
        R^{\text{cont}}_i \;\to\; R^{\text{cont}}_i
        + \bigl(n_i-n_i^{\text{prev}}\bigr)\frac{V_i}{\tau},
        \label{eq:dd-ptc}
        \end{equation}
        并让 $\tau$ 每成功一步 $\times10$（上限 $1$\U{s}）、每失败一步
        $\times0.1$；$\tau$ 足够大、Newton 进入收敛域后自动撤除，
        回到纯 Newton。
\end{enumerate}

\crefeq{eq:dd-ptc} 的关键在于\textbf{稳态残差本身不被改动}。时间项
只改变\emph{迭代方向}，Newton 循环始终在监视并判定\emph{物理}残差
$\mat{R}^{\text{phys}}$（实现上由 \texttt{aux['physical']} 单独携带）。
否则收敛判据会把"沿延拓路径走得怎么样"当成"稳态方程解得怎么样"，
两者不是一回事。

\subsection{平衡态锚点}
\label{sec:dd-equilibrium}

零偏压下的解不是"随便一个初值"，而是\emph{精确解}，
因此可以当作整个扫描的锚：先解一次非线性 Poisson
（$\psi$ 与 $n$ 在 Boltzmann 关系下自洽），把结果采用为接触储层参考：

\begin{equation}
\psi^{\text{src}}_j = \psi_{\text{eq}}(z_j, x_{\min}),
\qquad
\psi^{\text{drn}}_j = \psi_{\text{eq}}(z_j, x_{\max}).
\label{eq:dd-reservoirs}
\end{equation}

随后每个偏压点的接触条件由 \crefeq{eq:dd-reservoirs} 加端电压得到，
于是"接触面携带零电流"成为构造性的性质，而不是需要靠收敛去逼近的目标。

\subsection{偏压延拓的电位预移位}
\label{sec:dd-bias-shift}

把上一点的解直接当新初值有一个隐蔽的失败模式：漏极接触的电位跳了
$\Delta V_{ds}$，而内部还停在旧值上 —— 这个\emph{接触局部冲击}
足以让线搜索停滞，即便 Jacobian 是精确的。修法是按接触数据所强加的
同一个线性斜坡预移位内部电位：

\begin{equation}
\psi(x) \;\leftarrow\; \psi(x) + \Delta V_{ds}\cdot
\frac{x-x_{\min}}{x_{\max}-x_{\min}} .
\label{eq:dd-ramp}
\end{equation}

实测效果是把"停滞的解"变成"五步迭代的解"。此外，扫描本身按
\akey{max_step_V} 自动插入中间步（\emph{不}上报，只负责把解带过去）：
从 $0$ 直接跳到 $5$\U{V} 不是延拓，是带错初值的冷启动。

\subsection{端点电流与连续性门禁}
\label{sec:dd-continuity-gate}
\index{3method@方法!continuity gate@连续性门禁}

端点电流由接触面通量求和得到（进入器件为正）：

\begin{equation}
I_s \;=\; \sum_j J_{j,0}\,w_{z,j},
\qquad
I_d \;=\; -\sum_j J_{j,n_x-1}\,w_{z,j}.
\label{eq:dd-terminal-current}
\end{equation}

求和\emph{只遍历真正导通的面}（与残差使用同一掩码）。这一点不是细节：
若把势垒层的载流子面也算进去，端点电流会与求解器刚刚满足的
连续性平衡相差若干个数量级。

守恒性由门禁判定：

\begin{equation}
\frac{\abs{I_s+I_d}}{\max\bigl(\abs{I_d},\abs{I_s}\bigr)} \;\le\; \varepsilon_{\text{cont}}
\quad\text{当}\quad \abs{I_s+I_d} > I_{\text{floor}},
\label{eq:dd-gate}
\end{equation}

绝对地板 $I_{\text{floor}}$ 的存在是必要的：在平衡态附近端点电流是
机器零（实测 $2.7\times10^{-15}$\U{A/m}），此时\emph{相对}误差由浮点噪声
主导，对分母极小的比值设门槛只会制造假警报。低于地板者不设门槛。

\subsection{局部平衡闭合与辅助掩码}
\label{sec:dd-masks}

宽禁带势垒层内没有电流，其连续性方程对 $n$ 是线性齐次方程，
密度因此近乎\emph{自由}（Jacobian 的最小奇异值比最大低十个数量级）。
本模块用一个在物理上精确的替换消除这个自由模：在无电流行上以
\emph{局域平衡}闭合

\begin{equation}
n_i \;=\; \NCeff{i}\,\exp\!\left(\frac{\psi_i}{kT/q}\right)
\qquad\text{（取 } \varphi_n=0\text{）},
\label{eq:dd-local-equilibrium}
\end{equation}

取代连续性残差；行缩放随后把它变成一个\emph{相对密度误差}。
除这一掩码外，网格还携带另外四组掩码（欧姆接触、栅、背面、无栅表面，
以及载流子钉零的缓冲），它们的互斥性与对边界的覆盖由测试断言。

\section{验证与基准}
\label{sec:dd-verification}

验证纪律与 \crefchap{ch:sp1d} 一致：**先对闭式解，再谈不变量，
最后才看 IV 长什么样**。设计稿给出六个级别 T0--T5，
全部在 \fpath{tests/test_transport_dd.py} 中以断言形式固化。
制造解方法（MMS）的使用遵循其标准表述 \cite{roache2002,salari2000}，
验证与确认的总体框架见 \cite{oberkampf2010}。

\subsection{通量的闭式解与偏导}

\begin{veriftable}{Scharfetter--Gummel 通量：闭式解、反射恒等式与偏导}{tab:dd-flux}
\vrow{Bernoulli 函数对各区间}{$10^{-12}$ 相对}{\crefeq{eq:bernoulli} 的闭式解在宽区间内吻合}
\vrow{反射恒等式 $B(-x)=B(x)+x$}{恒等}{通过（含 $x\to0$ 与溢出区）}
\vrow{$B'(x)$ 对中心差分}{相对误差小}{通过（$10^{-12}$ 量级）}
\vrow{四个通量偏导对中心差分}{相对误差小}{通过（\emph{无}场依赖）}
\vrow{四个通量偏导对中心差分（启用 Caughey--Thomas）}{相对误差小}{
  通过（含 \crefeq{eq:dd-mobility-chain} 的 $\dd\mumu/\dd E$ 项）}
\vrow{势偏导反对称}{$\partial J/\partial\psi_L=-\partial J/\partial\psi_R$}{通过}
\vrow{密度偏导与密度无关}{\crefeq{eq:dd-flux-partials} 第一行}{通过（跨 $17$ 个数量级）}
\vrow{异质结台阶上的零通量}{$\abs{J}/\text{scale}<10^{-12}$}{通过}
\vrow{迁移率极限与导数}{低场 $\to\mumu_0$；$\dd\mumu/\dd E$ 对差分}{
  低场偏差 $<10^{-6}$；$E=0$ 处斜率 $<10^{-6}$}
\end{veriftable}

\subsection{T0：均匀电阻闭式解}

把网格退化为直通道、关断极化与肖特基之后，问题有闭式解
$J=qN_D\mumu V/L$。

\begin{veriftable}{T0\ \ 均匀电阻的闭式解（判据：相对误差 $<0.1\%$）}{tab:dd-t0}
\vrow{电流对闭式解 $J=qN_D\mumu V/L$}{$<10^{-3}$}{\textbf{$<10^{-6}$（超出判据三个数量级）}}
\vrow{两端电流一致}{一致}{$<10^{-9}$}
\vrow{电流与偏压严格正比}{比值 $=2$}{$2.0\pm10^{-6}$}
\end{veriftable}

\subsection{T1：平衡态不变量}

\begin{veriftable}{T1\ \ 零偏平衡态（判据：通量 $<10^{-12}$\U{A/m}）}{tab:dd-t1}
\vrow{平衡态端点电流}{$\le10^{-12}$\U{A/m}}{\textbf{$2.7\times10^{-15}$\U{A/m}（机器零）}}
\vrow{横向电流剖面}{$\le10^{-12}$\U{A/m}}{$\sim10^{-14}$\U{A/m}}
\vrow{Poisson 子问题平衡电荷平衡}{残差小}{$<10^{-10}$}
\vrow{带边算子与态密度台阶自洽}{重建 $=$ \akey{nc_eff}}{相对 $10^{-12}$}
\vnote{设计稿 T1 的另一半——与一维自洽 SP 的带边差 $<1$\U{meV}——\textbf{未做}，
      因为层界落在节点上的抹平使该判据在本离散下原理上不可达。
      这本身是一个诚实的结论，见 \S\ref{sec:dd-limits}。}
\end{veriftable}

\subsection{T2：制造解与二阶收敛}

Poisson 子问题用制造解检验：取 $\psi^{\star}=\sin(\pi z/L)$
（与 $x$ 无关，使横向项恒为零，从而\emph{隔离}竖直算子），
把解析源项加进去，加密四档。

\begin{veriftable}{T2\ \ Poisson 子问题的制造解收敛阶}{tab:dd-t2}
\vrow{竖直间距 $25$\U{nm}}{---}{$L_\infty$ 误差 $1.295\times10^{-2}$}
\vrow{$12.5$\U{nm}}{---}{$3.219\times10^{-3}$（观测阶 $2.008$）}
\vrow{$6.25$\U{nm}}{---}{$8.036\times10^{-4}$（观测阶 $2.002$）}
\vrow{$3.125$\U{nm}}{二阶收敛}{$2.008\times10^{-4}$（观测阶 $2.001$）}
\vnote{观测阶由相邻两档误差比的对数给出。三个值都贴在 $2$ 上，
      说明界面以外的格式确实是二阶的 —— 与界面处的一阶抹平并不矛盾：
      制造解在界面两侧都是光滑的。}
\end{veriftable}

\subsection{T3：端点电流连续性门禁}

\begin{veriftable}{T3\ \ 端点电流守恒（判据：相对误差 $<10^{-6}$，全部偏压点）}{tab:dd-t3}
\vrow{平衡态（$0$\U{V}）}{\crefeq{eq:dd-gate}}{低于绝对地板，按设计不设门槛}
\vrow{$V_{ds}\le1$\U{V} 的全部 $10$ 个偏压点}{$<10^{-6}$}{
  \textbf{$1\times10^{-13}$ 至 $1\times10^{-16}$}}
\vrow{耦合 Jacobian 对中心差分（逐行相对）}{相对误差小}{
  $\sim10^{-10}$（沿 Newton 路径）}
\end{veriftable}

\noindent
\textbf{"逐行相对"是判据的关键。} 本模块吃过一次大亏：$( \psi,\varphi_n)$
改 $(\psi,n)$ 时，Jacobian 里残留一行旧代码，把第二个未知量当成准费米势
代进了态密度公式 —— 内部密度与残差相差 \emph{17 个数量级}，
连续性行因此完全失去对 $\psi$ 的耦合，整个耦合矩阵在 Newton 最需要的
方向上近奇异。而当时的差分校验用\emph{列}最大值归一化，
坏掉的项绝对值很小，相对误差看着只有 $10^{-10}$。
改成\emph{逐行相对}判据后立刻暴露。这条教训值得写进任何后续模块：
\textbf{导数校验要在单元级（通量本身）做，且沿 Newton 路径做，
不能只在参考状态做一次。}

\subsection{T4--T5：延拓无关性与趋势}

\begin{veriftable}{T4\ \ 延拓无关性；T5\ \ 趋势与饱和}{tab:dd-t4-t5}
\vgroup{T4\ \ 初值/延拓无关性（判据：相对差 $<5\times10^{-3}$）}
\vrow{同一偏压点经两条不同阶梯到达}{相对差 $<5\times10^{-3}$}{通过（$0.05$ 与 $0.1$\U{V} 阶梯）}
\vgroup{T5\ \ 趋势合理性}
\vrow{栅压变负 $\Rightarrow$ 电流下降}{单调}{通过}
\vrow{迁移率加倍 $\Rightarrow$ 电流加倍}{比值 $=2$}{$2.0\pm0.1$}
\vrow{\textbf{漏电流饱和}}{高场斜率显著小于低场}{斜率比 $1.90\to1.75$}
\vnote{饱和判据必须用\emph{相对}对照。实测 $V_{ds}:0.25\to0.5\to1.0$\U{V} 时
      $I_d:146\to278\to478$\U{A/m}；同器件在\emph{常数}迁移率下斜率比为
      $2.00\to0.92$ —— 栅控器件本身就带约 $8\%$ 的次线性
      （栅-沟道电位改变二维电子气），所以不能指望对照组是欧姆的。
      测试断言的是"场依赖使斜率比明显更小"这一相对判据。}
\end{veriftable}

\subsection{端到端：演示 Deck}

\begin{veriftable}{演示 Deck 的端到端结果（\fpath{decks/demo_hemt_dd.json}）}{tab:dd-demo}
\vrow{运行状态}{success}{通过（$10$ 个偏压点全部收敛）}
\vrow{墙钟时间}{---}{约 $50$\U{s}（\akey{max_step_V} $=0.05$\U{V}）}
\vrow{二维电子气面密度与极化上界}{\crefeq{eq:ns-bound}}{
  比 $\sigeff/q=9.6\times10^{12}$\U{cm^{-2}} 高约 $\mathbf{1.7}$ 倍}
\vrow{网格掩码互斥且覆盖边界}{互斥 + 覆盖}{通过}
\vnote{1.7 倍这一偏差可归因于层界落在节点上的抹平
      （\crefchap{ch:mesh} 与 \S\ref{sec:dd-discretisation}），
      设计稿已预见，此处首次量化。$0.25$\U{V} 阶梯只需约 $9$\U{s}，
      但末端点不收敛 —— 时间是拿收敛性换来的。}
\end{veriftable}

\section{误差来源与适用边界}
\label{sec:dd-limits}
\index{1concept@概念!validity limits@适用边界}

\limitsheading

\begin{bounds}
  \item \textbf{迁移率是占位参数。} \akey{mu0_cm2_Vs}（$1500$\U{cm^2/(V\cdot s)}）
        与 \akey{vsat_cm_s}（$2\times10^{7}$\U{cm/s}）均未经标定。
        L2 基准之前，\emph{绝对数值不可信}；可用的是趋势、
        守恒性与收敛阶。

  \item \textbf{只覆盖到约 $1$\U{V} 的漏压。} 速度饱和已实现且默认启用，
        但 $V_{ds}\ge2$\U{V} 时 Newton 仍不收敛（欠账）。
        这是模型边界而非调参问题。

  \item \textbf{短栅的速度过冲与弹道成分缺席。}
        $L_g=150$\U{nm} 已在漂移扩散假设开始失效的尺度上。
        方案 §15 的对策（流体动力学模型插件 + 经验修正 + 报告标注）
        原样适用。

  \item \textbf{界面被一个单元宽度抹平。} 量化后果即 $1.7$ 倍的面密度偏差。
        设计稿 T1 中"与一维 SP 的带边差 $<1$\U{meV}"这条判据
        在本离散下\emph{原理上不可达} —— 这是诚实结论，不是待办事项。

  \item \textbf{只实现 Boltzmann 统计。} \akey{statistics} 接受其他取值时
        显式报错，因为简并修正必须同时改通量；在 \crefeq{eq:dd-poisson}
        意义上的高密度沟道中，这会使面密度与电流有系统性偏差。

  \item \textbf{缓冲层是抽象而非器件。} "中性介质 + 载流子钉零 + 零通量壁"
        是单极 HEMT 的标准处理，但它不是真实的半绝缘缓冲：
        深能级、俘获、背栅效应一概没有。

  \item \textbf{无栅电流。} $I_g$ 由构造为零，而非由求解得到。
        因此一切与栅泄漏相关的现象（栅漏电、击穿）不在模型内。

  \item \textbf{无陷阱、无自热、无瞬态。} 电流坍塌、gate-lag、
        drain-lag、晶格温升一概没有；RF 大电流点的预测因此偏乐观。

  \item \textbf{无栅表面钉扎是工程近似。} 真实的无栅表面由施主态钉在
        特定能级；本模块把它钉在 Schottky 参考电位上。
        它同时承担"给势垒电位一个锚"的数值职责 ——
        去掉它，势垒电荷近似为零时 Poisson 退化为纯 Laplace，
        势垒电位存在近乎自由的模式。

  \item \textbf{PTC 与阻尼的参数是工程值。} $\tau_0$（$10^{-12}$\U{s}）
        与减半次数不是物理量；换网格尺度可能需要重调。
        测试只锁\emph{行为}，不锁这些参数。
\end{bounds}

\section{接口契约与小结}
\label{sec:dd-interface}
\index{5artifact@数据通道!iv}\index{5artifact@数据通道!current}

{\small
\begin{longtable}{@{}p{0.06\textwidth} p{0.24\textwidth} p{0.60\textwidth}@{}}
\toprule
\textbf{方向} & \textbf{键} & \textbf{本模块使用/提供的字段}\\
\midrule
\endfirsthead
\toprule
\textbf{方向} & \textbf{键} & \textbf{本模块使用/提供的字段}\\
\midrule
\endhead
\bottomrule
\endfoot
入 & \akey{mesh} & \akey{x_nm}, \akey{z_nm}, \akey{nx}, \akey{nz}, \akey{layers[]}（含 \akey{node_slice}）\\
入 & \akey{structure} & \akey{layers[].{role, material, x_al, doping_cm3, doping_kind}},
  \akey{gate.barrier_eV}\\
入 & \akey{materials} & 介电常数、有效质量、带隙（含合金插值）\\
入 & \akey{polarization} & $\sigeff$、势垒厚度、$\Delta E_g$（$\dEc$ 的优先来源）\\
入 & \akey{quantum}（可选） & \akey{potential_eV} 或 \akey{psi_V}，仅作\emph{初值}\\
出 & \akey{electrostatics} & \akey{psi_V}、极值、坐标轴、\akey{gate_reference_V}、偏压\\
出 & \akey{carriers} & \akey{density_m3}、极值、\akey{ns_cm2}、\akey{statistics}、迁移率\\
出 & \akey{current} & 逐点记录（收敛、迭代数、残差、$I_d$/$I_s$/$I_g$、
  连续性误差与门禁结果）、最差连续性误差、失败清单\\
出 & \akey{iv} & \akey{vds_V}/\akey{vgs_V}/\akey{id_A_per_m} 曲线，
  单位为 \U{A/m}（每单位器件宽度）\\
出 & \akey{ns} & 最后偏压点的面密度（\U{cm^{-2}}）\\
\end{longtable}
}

\noindent
契约上有两处刻意的设计：\textbf{场数组只记最后一点}，而\emph{每个}
偏压点的摘要都在 \akey{current} 里 —— 二维场数组随偏压点线性增长，
全存会撑爆上下文，而"哪一点没收敛、为什么"才是诊断真正需要的。
另一处是 \akey{iv} 的逐点完整记录（设计稿只要求曲线，实现多加了
\akey{points}）：门禁失败必须能定位到\emph{具体偏压点}，只给曲线不够用。

\subsection*{小结}

本章的结论可以概括为三点。

其一，\textbf{未知量的选择是一个物理判断，不是数值偏好。}
取 $n$ 而非 $\varphi_n$ 作未知量，使宽禁带势垒与耗尽区的连续性方程
保持良态（$\partial J/\partial n$ 与 $n$ 无关）—— 代价是动态范围极大，
由\emph{固定物理参考量}的行缩放处理，而不是自适应对角均衡。

其二，\textbf{三层保险解耦了三件不同的事}：延拓负责给出好的初值，
线搜索负责保证每一步都下降，PTC 负责在 Newton 暂时失效时改变
\emph{迭代矩阵}而不改变\emph{被解的方程}。稳态残差始终是判据本身，
因此"沿延拓路径的表现"永远不会被误当成"稳态方程的解"。

其三，\textbf{验证的层级必须先于"像不像器件"。}
T0--T5 的顺序不是形式：闭式解先证明\textbf{离散是对的}，
不变量再证明\textbf{守恒是对的}，收敛阶证明\textbf{格式是二阶的}，
只有在这些之后，IV 曲线上的数字才值得被引用 ——
而在迁移率标定之前，那些数字仍然只是\emph{趋势}。

\cleardoublepage
